\section{Passages with two logarithmic scales}\label{sec:mixed}

Let $L,W$ be positive inputs tending to infinity, and write $c=W/L$.
An occurrence of $c$ in a layer or in the homogeneous part of the
equation always denotes this exact ratio. Other bounded parameters of
the field are independent ordinary parameters, even if the original
matching problem subsequently ties them to a function of this ratio.
We consider the following scalar equations on $0\leq\theta\leq L$:
\begin{align}
 v'&=-cv+g(X,Y,v),&g(0,0,v)&=0,\label{eq:mixed-I}\\
 z'&=G(X,Y,z).&&\label{eq:mixed-II}
\end{align}
In Equation~\eqref{eq:mixed-I}, the left and right layers are
\[
 X_\ell=e_\ell e^{-a_\ell\theta},\qquad
 Y_\ell=d_\ell e^{-b_\ell(L-\theta)},\qquad a_\ell,b_\ell\in\mathbb Q_{>0}.
\]
In Equation~\eqref{eq:mixed-II}, their rates have the form
$\alpha_\ell+\beta_\ell c$, with rational $\alpha_\ell,\beta_\ell$ and
$c\to c_0<\infty$. If $c_0>0$, all limiting rates are positive and
$G(0,0,z)=0$. If $c_0=0$, a rate either has $\alpha_\ell>0$, or has
$\alpha_\ell=0$, $\beta_\ell>0$; in this case $G$ vanishes when all
layers of the first kind are set to zero. We call those layers fast.
Thus every field monomial contains a fast layer. For positive $c_0$
every layer can be designated fast.

For prescribed input $v(0)=v_{\rm in}$ or $z(0)=z_{\rm in}$, the
output is the state at $\theta=L$. The clocks, layer amplitudes, input
state, and ordinary field parameters are independent arguments of this
output; the occurrences of $c$ in the layers and homogeneous generator
are evaluated by the exact formula $W/L$. Endpoint and coefficient ties
are imposed later on these arguments.

Fields are analytic and bounded on larger complex boxes than those
used for evaluation. The amplitudes have sufficiently small fixed
modulus; in Equation~\eqref{eq:mixed-I} the input state is also small,
whereas in Equation~\eqref{eq:mixed-II} it may range over a compact
subbox of the state domain. All parameter statements below hold on
smaller ordinary complex neighborhoods, uniformly on their closures.

\begin{theorem}[Mixed passage packets]\label{thm:mixed-packets}
Suppose that the primary inputs are affine combinations of flag scales
with constant coefficients and bounded analytic parts, and that all
large basis nodes in those combinations have dependent logarithms.
For an unbalanced limit, put $P=\max(L,W)$ and $R=\min(L,W)$ on the
reference sequence, and include a primary input $C_{\rm aux}\to\infty$
tied there to $P/R$. Assume the saturated flag and controlled tests of
Sections~\ref{sec:flags} and~\ref{sec:calculus}. Then the outputs of
Equations~\eqref{eq:mixed-I} and~\eqref{eq:mixed-II} satisfy the
primitive hypotheses of Assumption~\ref{ass:calculus-primitive},
including fixed independent derivatives, simultaneous safe columns,
and refinement invariance of terminal analytic germs. They are
retestable for elimination in the sense of
Definition~\ref{def:elimination-retestable}.

The amplitude and input smallness is independent of expansion depth.
For a family of positive limiting slopes it can also be chosen
independently of the slope, provided the ordinary boxes and field
bounds are common and every field monomial contains a designated
layer whose limiting rate has a common positive lower bound. The
other rates need only be positive at each particular limiting slope.
Constants and starting points for a fixed packet label may depend on
that slope. The two unbalanced cases have common smallness for each
fixed rate list.
\end{theorem}

The three ratio regimes use different coefficient constructions.
For comparable scales write $H=W-c_0L$; for an unbalanced passage
write $C=P/R$, with $P,R$ as in the theorem.  Here $P$ is the
larger primary action, whereas in the preceding two sections it
denoted a small transfer charge.
\begin{center}
\small
\begin{tabular}{@{}>{\raggedright\arraybackslash}p{0.20\linewidth}>{\raggedright\arraybackslash}p{0.28\linewidth}>{\raggedright\arraybackslash}p{0.44\linewidth}@{}}
\toprule
Ratio regime & Construction & Intermediate coefficient data\\
\midrule
$c\to c_0\in(0,\infty)$ & Integration trees and residues
 & Analytic cluster functions when $H$ is bounded; separated exact
shifts with explicit power prefactors and analytic coefficient jets
when $H$ is unbounded\\[3pt]
$c\to0$ & Integration trees, then cutoff retirement
 & Meromorphic ratio coefficients, then Gevrey series in $1/C$\\[3pt]
$c\to\infty$ (model~\eqref{eq:mixed-I}) & Compatible input and output charts
 & Meromorphic chart coefficients, then Gevrey series in $1/C$\\
\bottomrule
\end{tabular}
\end{center}
Each route is local to its chosen limiting regime.  The ratio in
the layers and homogeneous generator remains the exact $W/L$;
other ordinary field parameters remain independent until their
argument ties are imposed.

We first prove the angle estimates used by all the constructions.
The pole estimates enter with the unbalanced coefficient formulas.
In this section a constant attached to a fixed charge,
coefficient, or derivative order need not be uniform as that finite
request increases.

\subsection{Primary combinations and ratio geometry}

Write the flag-path parameter as $s=u+i v_{\rm ang}$, distinguishing
its imaginary coordinate from the state variable $v$.

\begin{lemma}[Ratio geometry]\label{lem:mixed-ratio}
If a positive primary combination $J$ has leading index $p$, then
\begin{equation}\label{eq:mixed-primary}
 J=A_JZ_p(1+\Theta_J),\qquad A_J>0,
\end{equation}
where $\Theta_J$ is a finite sum of strict-small monomials, of first
index greater than $p$, times bounded analytic factors. On every
usable band $j\leq p$, including its required side,
\begin{equation}\label{eq:mixed-primary-real}
 |J|\asymp |Z_p|,\qquad
 \Re J\asymp |Z_p|\cos(\Im f_p)\gg(\log|J|)^q
 \quad(q\text{ fixed}).
\end{equation}
For an unbalanced passage, let $i,n,k$ be the leading indices of
$P,R,C_{\rm aux}$, respectively. Then $n,k>i$, possibly $n=k$, and
there is an exact identity of flag expressions
\begin{equation}\label{eq:mixed-log-tie}
 f_i=f_n+f_k+b_*,\qquad
 C:=\frac PR=\frac{A_P}{A_R}e^{b_*}Z_k
                  \frac{1+\Theta_P}{1+\Theta_R}.
\end{equation}
Here $b_*$ is bounded background data, $\Lambda_i=\Lambda_n+\Lambda_k+O(1)$,
and $\sigma(i)>\min(n,k)$.

Let $C_\bullet$ denote the same ratio formula with any already
dispersed small arguments set to zero in its units. While it is used
meromorphically, on bands of index at most $k$,
\begin{align}
 |C_\bullet|&\asymp |Z_k|,&
 \arg C_\bullet&=(1+o(1))v_{\rm ang}\Lambda_k,\label{eq:mixed-ratio-angle}\\
 |\Im C_\bullet|&\gtrsim |Z_k|\min(1,|v_{\rm ang}|\Lambda_k).&&\label{eq:mixed-ratio-im}
\end{align}
The lifted argument is in $[-\pi/2-o(1),\pi/2+o(1)]$, and is
transverse to the real axis at the $k$ fringe. If
$|v_{\rm ang}|\Lambda_k\ll1$, then
\begin{equation}\label{eq:mixed-real-shift}
 \Re C_\bullet=C_\bullet(u)
             \bigl(1+O((v_{\rm ang}\Lambda_k)^2)\bigr).
\end{equation}
Through the $i$ band the actual positive actions $P,R$ have increasing
arguments, $0<\arg R<\arg P<\pi/2$ on the positive side. At its fringe,
\begin{equation}\label{eq:mixed-cos-gap}
 \frac{\cos\arg R}{\cos\arg P}\longrightarrow\infty,
 \qquad |\sin(\arg P-\arg R)|\gtrsim\frac{\Lambda_k}{\Lambda_i}.
\end{equation}
The negative side has the reflected assertions.
\end{lemma}

\begin{proof}
Divide the affine formula for $J$ by its leading term. Ratios of
primary nodes, and the reciprocal of the leading node multiplying a
bounded term, have exact logarithms in later nodes. Their first
nonzero exponents are positive, which proves
Equation~\eqref{eq:mixed-primary}. On the real ray their moduli tend
to zero. On a $j$ band the flag angle estimates control their real
parts as well: if $l>p$ and the $p$ cosine is small, the possible
relative cosine gain is at most a constant times
\[
 1+\frac{\Lambda_p-\Lambda_l}
 {\Lambda_j\cos(\Im f_j)+\Lambda_j-\Lambda_p}.
\]
The logarithmic-gap derivative estimate absorbs this factor in
$|Z_l/Z_p|$. If the $p$ cosine is small, $\Lambda_j=O(\Lambda_p)$;
the polynomial side buffer therefore remains a polynomial buffer in
$f_p(u)$. This proves Equation~\eqref{eq:mixed-primary-real}, including
the bounded affine part. The same estimates after differentiation
show that the argument of the unit in
Equation~\eqref{eq:mixed-primary} is $o(|v_{\rm ang}|)$.

On the reference sequence the assumed tie makes
$f_i-f_n-f_k$ bounded. Each of these logarithms has an exact
triangular affine expression. A nonzero coefficient of its largest
remaining large scale would make their difference unbounded, so
every large coefficient cancels. Their difference is exactly a
bounded background combination $b_*$. Neither $f_n$ nor $f_k$ uses
a node at or before $\min(n,k)$, proving the assertion about
$\sigma(i)$. Differentiation gives the assertion about $\Lambda_i$.
This argument uses the auxiliary tie on the reference sequence; it
does not impose that tie at independent ambient test points.

The triangular angular formulas applied to the surviving expressions
in Equation~\eqref{eq:mixed-log-tie} prove
Equations~\eqref{eq:mixed-ratio-angle} and~\eqref{eq:mixed-ratio-im}.
Deleting an already dispersed argument preserves the proof: every
remaining strict-small term of first index $l\geq j$ is
$O(e^{-a\Re Z_l})$, and its fixed angular derivatives lose only
$\exp(O(f_l(u)))$. For controlled ordinary backgrounds use their
real Taylor jets before this comparison. Downward differentiation
bounds the sensitivity of $f_l$ by
$f_l(u)\operatorname{polylog}f_l(u)$ on the relevant angular scale;
a sufficiently high fixed common jet thus preserves any specified
polynomial angular margin. Real symmetry removes the linear term
from the real part, giving Equation~\eqref{eq:mixed-real-shift}.
In particular $|\Im C_\bullet|<1$ implies
$\Re C_\bullet-C_\bullet(u)=O(|C_\bullet(u)|^{-1})=o(1)$.

Finally $\Lambda_i-\Lambda_n=\Lambda_k+O(1)$, with a diverging
positive gap. The good-to-side angular range gives
Equation~\eqref{eq:mixed-cos-gap}. The flag side dominance also makes
$\Re R/\Re P$ smaller than any prescribed fixed positive constant
after its dominance constant is chosen. All comparisons have slack
under the smaller neighborhoods subsequently used.
\end{proof}

\subsection{The actual analytic passage and integration trees}

\begin{lemma}[Continuation in the leading band]\label{lem:mixed-actual}
Let $i$ be the leading index of the larger primary action, or the
common leading index of $L,W$ when they are comparable. The actual
passage is bounded and holomorphic throughout band $i$, with the
parameter and state room
needed for fixed independent derivatives.
\end{lemma}

\begin{proof}
For comparable scales put $H=W-c_0L$ and $\delta=H/L$; $H$ is
bounded ordinary data or a signed primary combination with later
leading index. On the $i$ band, $\delta\to0$. At its fringe put
$\rho=\Re L/|L|$ and denote the side dominance constant by
$\mathcal A$. The exact identity for $\Im(H/L)$ gives
\begin{equation}\label{eq:mixed-comparable-detuning}
 \frac{|\Im\delta|}{\rho}
 \leq\frac{|H|}{|L|}+\frac{|\Re H|}{\Re L}
 \leq\frac{C_H}{\mathcal A}+o(1).
\end{equation}
Indeed $H$ has only later affine nodes and bounded terms, whereas
the side inequality makes $\Re Z_i$ dominate $\mathcal A$ times
the sum of the later real parts. Increasing $\mathcal A$ once
also makes $\Re L\geq (A_L/2)\Re Z_i$. The first term tends
to zero by the modulus hierarchy, and the bounded terms are
negligible by the polynomial buffer. These estimates are uniform
for the controlled families of
Definition~\ref{def:calculus-controlled-family}; $C_H$ depends on
the fixed affine formulas and common ordinary bounds, not on the
moving anchor. The last quantity need not tend to zero for fixed
$\mathcal A$.

Fix $0<\kappa\leq1/4$ before choosing any expansion budget.
For the finite rate list let $K_0$ bound each generator's absolute
detuning coefficient divided by its positive limiting rate,
including the homogeneous generator when present. Choose
$\mathcal A$ so that $2K_0C_H/\mathcal A<\kappa/4$, and then
impose the numerical far thresholds
$K_0|\delta|<1/4$ and the corresponding small bound on the
$o(1)$ in Equation~\eqref{eq:mixed-comparable-detuning}.
Every fixed positive limiting action then retains a fixed
positive fraction of its real action: for a generator with
limiting rate $q_h>0$ and detuning coefficient $\beta_h$,
\[
 \Re\bigl[(q_h+\beta_h\delta)L\bigr]
 \geq q_h\Re L-|\beta_h|\,|\Re H|
 \geq \tfrac12 q_h\Re L.
\]
For any sum of generator
costs $\nu+B\delta$ one also has
$|\arg(\nu+B\delta)|\leq\kappa\rho/2$.
The cone parameter is fixed independently of the later budget;
the side constant and thresholds may depend on the limiting slope.

For comparable scales and for $c\to0$, integrate on the polygon
\[
 0,\ S,\ L-S,\ L,\qquad S=K\log|L|>0.
\]
Fast real rates are bounded below on the horizontal ends. Along
the middle every left and right fast profile has an exponential
margin from the corresponding end; choosing $K$ large pays its
length. The integral of the sum of the absolute fast profiles is
therefore at most a constant times the amplitude size, uniformly in
$L$. Pure slow profiles stay bounded. Homogeneous $c$ actions, when
present, have increasing real part: $\Re c\geq0$ and
$|L|\cos\arg W\gg S$. Signed corrections to fast rates are
absorbed by $\Re W/\Re L\ll1$ in the small-slope case, and by the
preceding $\delta$ estimate in the comparable case.

For $c\to\infty$, use the straight segment when $\cos\arg L$
has a positive fixed lower bound, and the same polygon otherwise.
In the latter case, including a switching overlap,
$\Lambda_n\gtrsim\Lambda_i$, so the real path estimates give
$f_i\leq\operatorname{poly}(f_n)$. The buffer for $W$ then gives
$|L|\cos\arg W\gg S$, as required for monotone damping.
The integral of $|g_v|$ is small in all cases, because each field
monomial has a fast factor; the stable propagator has modulus at
most one. Picard iteration on these paths is a contraction in a
fixed state box, after the amplitudes are reduced once. The same
estimate applied to the forcing keeps the state inside a smaller
box. Ordinary parameter Cauchy margins give the claimed room.

Local deformations of the polygons preserve their estimates and
give holomorphic dependence. On overlaps uniqueness identifies the
solutions, including an interpolation to straight segments in the
switching region. Their union is the continuation of the real
passage. Before index $i$ use this actual function with zero
prefixes.
\end{proof}

For comparable or small slopes, expand the integral equation by
analytic substitution. A planar rooted integration tree has a field
monomial at each vertex, and therefore at least one fast factor
there. In model~\eqref{eq:mixed-I}, homogeneous input leaves and the
linear edge propagators are included. The free term is respectively
$e^{-W}v_{\rm in}$ or $z_{\rm in}$. Let $r$ be the number of
vertices and $M_0$ the total layer degree plus homogeneous-input
leaf count. Then $M_0\geq r$, and the sum of vertex arities is
$O(M_0)$. Cauchy bounds for the field, the finite number of layer
generators, and the exponential count of planar trees give total
weights bounded by $\eta^{M_0}$ after a fixed exponential factor is
absorbed into a sufficiently small $\eta$. Integrating one fast
factor per vertex on the polygons of Lemma~\ref{lem:mixed-actual}
proves absolute convergence and the equation by substitution.
In model~\eqref{eq:mixed-II} expand about the bounded input, using
its fixed state margin, so it does not have to be small.

Each tree integration domain is the union of at most $r!$ ordered
simplexes. On one such simplex the $r+1$ successive gaps have costs
$C_j=A_j+cB_j$. Its coefficient, after removing the monomial weight,
is the convolution of the $r+1$ exponentials, and hence
\begin{equation}\label{eq:mixed-residue}
 \sum\operatorname*{Res}_{s=-C_j}
       e^{sL}\prod_{j=0}^r(s+C_j)^{-1}.
\end{equation}
For distinct costs this follows by partial fractions; the integral
and the total residue are entire in all costs, which proves the
identity also at collisions. Straight paths can be used for this
identity because the simplex integrands are entire.

\begin{example}[Collision of two costs]\label{ex:mixed-cost-collision}
For one event the convolution is
\[
 \int_0^L e^{-C_0t}e^{-C_1(L-t)}\,\dd t
 =\begin{cases}
 \displaystyle\frac{e^{-C_0L}-e^{-C_1L}}{C_1-C_0},&C_0\ne C_1,\\[6pt]
 Le^{-C_0L},&C_0=C_1.
 \end{cases}
\]
The separate residue terms have a denominator at a cost collision,
whereas their sum extends analytically there.  The cluster integral
below keeps this cancellation before the corresponding charge is
dispersed.
\end{example}

\begin{lemma}[Ordered cost and contour bounds]\label{lem:mixed-tree}
For comparable slopes put $\nu_j=A_j+c_0B_j$. These values belong
to a finitely generated positive-weight semigroup and satisfy
\begin{equation}\label{eq:mixed-comparable-costs}
 |B_j|\leq K_0\nu_j,\qquad
 \nu_j+\nu_{j'}\geq a|j-j'|.
\end{equation}
For small slopes, $A_j$ is in a nonnegative rational lattice, $B_j$
in a signed rational lattice, and
\begin{equation}\label{eq:mixed-small-costs}
 A_j+A_{j'}\geq a|j-j'|,\qquad
 B_j\geq-K_0A_j,\qquad |B_j|\leq K_0M_0.
\end{equation}
The $l$th main cost in increasing order is at least $a(l-1)/2$.
For every large fixed $N$, the aggregate high-residue sum at the
leading fringe has modulus at most
\begin{equation}\label{eq:mixed-high-residues}
 \frac{K_N K_*^{r+1}}{r!}\,e^{-bN\Re L},
\end{equation}
where $b>0$ and the exponential base $K_*$ are independent of $N$.
In the comparable case ``high'' means main costs above a threshold
in $[2N,3N]$. In the small-slope case it means projected costs
$p_j=A_j+\mu B_j$ above such a threshold, where
$\mu=\Re W/\Re L$.
\end{lemma}

\begin{proof}
Every event between gaps $j$ and $j'$ supplies a fast rate to one of
their costs: a left layer contributes before its event and a right
layer after its event. Nonnegative homogeneous damping contributes
additional cost. This proves the pairwise lower bounds. For
comparable slopes each generator has a positive limiting weight,
so its $B$ contribution is bounded by a fixed multiple of that
weight. For small slopes the negative $B$ contributions occur only
in fast rates and are bounded by their $A$ contribution; pure slow
rates and homogeneous damping have nonnegative $B$. The total
degree bounds the absolute sum of $B$ contributions. If $l$ costs
are at most $t$, their first and last original indices differ by
at least $l-1$, whence $2t\geq a(l-1)$. This proves the sorted bound.

We give the contour estimate since it keeps the smallness of the
analytic inputs independent of $N$. Reflecting if necessary, take
the positive fringe and write $\rho=\Re L/|L|\to0$;
$\Re L\gg\log(1/\rho)$. Fix the same $0<\kappa\leq1/4$ as in
Lemma~\ref{lem:mixed-actual}. In the comparable case
Equation~\eqref{eq:mixed-comparable-detuning}, with its specified
choice of side constant, puts every pole within angle
$\kappa\rho/2$ of the negative real ray. Choose
$N'\in[2N,3N]$ separated from the finitely many base semigroup
values there. Enclose the high poles by a cap $\Re s=-N'$ and
rays of angles $\pi\pm\kappa\rho$. Since $\kappa<1$, both rays
have a fixed positive exponential-decay margin. Distances on the cap have a positive
fixed lower bound for the finitely many low costs. On the rays
at modulus $x$, at most $K(x+1)$ sorted indices require the
distance bound $b\rho x$; all others have distances bounded below
by a fixed multiple of their sorted index. The factorial of those
initial indices is bounded by $K^m x^m$ if their number is $m$.
Consequently
\[
 r!\prod_j|s+C_j|^{-1}
 \leq K_*^{r+1}(K_*/\rho)^{K_*(x+1)}.
\]
On the cap the last factor can be replaced by $K_N\rho^{-K_N}$.
Along the rays $\Re(sL)\leq-bx\Re L$, so the latter factors are
absorbed by half the decay sufficiently far out. Integration over
the rays and cap proves Equation~\eqref{eq:mixed-high-residues}.
The same estimate makes the closing arcs tend to zero; all poles
of each finite product are included before the limit is taken.

For small slopes, side dominance makes $\mu$ small, whereas
Equation~\eqref{eq:mixed-cos-gap} gives $\mu/|c|\to\infty$.
Thus $p_j\geq A_j/2$. The number of projected costs below $3N+1$
is bounded in terms of $N$, so for each simplex there is a threshold
$N'\in[2N,3N]$ a fixed distance from all these projected costs.
The poles lie in the cone between angles $\pi/2$ and
$\pi+\kappa\rho/2$: when $B_j<0$ use
$|c|\sin(\arg L-\arg W)\leq\mu\rho$ and
$|B_j|\leq K_0A_j$. Use the cap
$\Re(sL)=-N'\Re L$ and rays of angles
$\pi+\kappa\rho$ and $\pi/3$. The cap has length and modulus
$O(N')$, distances at least $b_N\rho$ for its finitely many small
costs, and factorial distances for the large $A_j$'s. On the first
ray the preceding estimate applies. On the second ray, starting
at modulus comparable to $N'\rho$, the exponential decays as
$e^{-bx|L|}$ and angular separation gives distances at least $bx$.
For $x<1$ only a fixed number of sorted indices require that bound,
so a fixed inverse power of $\rho$ suffices; for $x\geq1$ use the
same factorial argument. This proves the stated estimate also in
this case. All low projected costs have $A_j\leq2N'$ and, since
$\mu/|c|\to\infty$, $B_j<\eps |Z_k|$ for every fixed positive
$\eps$ sufficiently far out.
\end{proof}

\subsection{Comparable slopes and low residues}

Fix a main comparable charge $\nu$ and put $J=\{j:\nu_j=\nu\}$.
Lemma~\ref{lem:mixed-tree} bounds $|J|$ in terms of $\nu$.
The factor
\begin{equation}\label{eq:mixed-Knu}
 K_\nu(t,\delta)=
 \prod_{j\notin J}(t+\nu_j-\nu+B_j\delta)^{-1}
\end{equation}
is analytic on a fixed small $(t,\delta)$ bidisk, with bound
$K^{(\nu)}K_*^{r+1}/r!$. Indeed finitely many low distinct values
have a positive separation, and every sufficiently high value
pays a fixed multiple of its sorted index. This also bounds any
fixed jet after a disk shrink.

If $H=W-c_0L$ is bounded, the whole cluster contributes
$e^{-\nu L}L^{|J|-1}$ times
\begin{equation}\label{eq:mixed-bounded-H}
 \frac{1}{2\pi i}\int_\Gamma
 \frac{e^zK_\nu(z/L,\delta)}
 {\prod_{j\in J}(z+B_jH)}\,\dd z,
\end{equation}
where a fixed circle encloses these finitely many bounded poles
uniformly on a smaller ordinary neighborhood. The integral is an
analytic function of its small and ordinary arguments, with the
same factorial bound. If $H$ is unbounded, split the cluster into
its distinct $B$ values, a finite set for fixed $\nu$. Direct
residue differentiation gives exact shifts $e^{-\nu L-BH}$,
boundedly many powers of $L$ and $\delta^{-1}$, and analytic
values or jets of Equation~\eqref{eq:mixed-Knu}.

The $r!$ possible linear extensions are canceled by the factorial
gain, and the remaining $K_*^{r+1}$ is paid by the original
geometric weights. Thus summing these coefficients over trees is
normally convergent on fixed smaller argument boxes. The
$K^{(\nu)}$ and disk sizes may depend on the charge but occur only
once per fixed low-charge expression, rather than once per vertex.
The error in Lemma~\ref{lem:mixed-tree}, summed in the same norm,
pays arbitrary depth at the $i$ fringe by increasing $N$.

Every remaining power or unit, including $H/L$, is an exact
monomial times an analytic unit in strict-small arguments. A
bounded $H$ has only ordinary terms. Split the affine exponential
shifts exactly and Taylor-disperse each small argument at its first
index, using the calculus of Section~\ref{sec:calculus}.
Taylor remainders pay arbitrary finite depth after fixed prefactor
losses; the support is iterated left-finite. This constructs the
coefficient functions and transitions in the comparable case. Their
differentiated bounds and terminal-germ invariance are collected with
the unbalanced cases below.

Under the common fast-rate hypothesis in
Theorem~\ref{thm:mixed-packets}, use that fast layer at each event
in the path integral and sorted-index estimate. The constant $a$
and the exponential base $K_*$ can then be fixed uniformly. Shrink
the detuning neighborhood, separately for each slope, to keep the
large-cost distances above a common multiple of their base values.
Only the low-charge constants and radii change. The same original
amplitude choice therefore works throughout the asserted family.
More explicitly, fix the cone parameter $\kappa\leq1/4$ and the
common designated-rate lower bound $a_*>0$ first. The slope-specific
$K_0$ is paid solely by the side constant $\mathcal A$ and the
numerical thresholds in Equation~\eqref{eq:mixed-comparable-detuning}.
After those choices, every designated profile has horizontal decay
rate at least $a_*/2$, and its real middle action has the same fixed
fractional margin. Its end integrals are bounded by $2/a_*$, while
a common choice in $S=K\log|L|$ pays the middle length. All other
profiles are bounded by their amplitudes. Thus the integral and
tree norm constants are common. On the contour, at most
$m\leq C_{a_*}(x+1)$ initial sorted indices need the weak distance
bound, and
\[
 r!\prod_j|s+C_j|^{-1}
 \leq K^{r+1}(m-1)!(\rho x)^{-m}
 \leq K_*^{r+1}(K_*/\rho)^{C_{a_*}(x+1)}
\]
for $x\geq1$, with common $K_*$. Neither $\mathcal A$ nor $K_0$
appears once per vertex. A later budget $N$ affects only its finite
cap/low-charge constant and thresholds; increasing side dominance
for that request improves the already fixed cone margin.

\subsection{Small damping and retirement of its meromorphic coefficients}

Here $P=L$, $R=W$, $C=L/W$. The low-charge residue formulas are
meromorphic in $C$. We retain these formulas up to the ratio's fringe,
where they can be replaced by Gevrey series in $1/C$.

\subsubsection{Retained low-charge residues}

For each fixed main charge $A$, retain
the residues with $B$ below a cutoff comparable to
$\eps|Z_k|$, with $\eps>0$ fixed and small. Freeze this cutoff
locally in a fixed-factor modulus bin. The multiplicity $m_{AB}$
of an identical pair $(A,B)$ is bounded in terms of $A$.
Expanding the other factors to order $m_{AB}-1$ yields terms
\begin{equation}\label{eq:mixed-small-residue}
 e^{-AL-BW}L^d C^t M_{ABdt}(1/C),
\end{equation}
where $d,t\geq0$ range over finite sets for fixed $A$. All factors
with equal $A$ and unequal $B$ have a fixed lattice gap in their
$B$ difference; their powers of $C$ have been included in $C^t$.
The remaining functions are fixed jets of products with factors
\begin{equation}\label{eq:mixed-small-denominator}
 \bigl(A_j-A+(B_j-B)/C\bigr)^{-1},\qquad A_j\ne A.
\end{equation}

For $B$ in the cutoff, every sufficiently large $A_j$ pays a
constant multiple of $A_j$: use $B_j\geq-K_0A_j$, bounded
$|B/C_\bullet|$, and the fact that $1/C_\bullet$ stays away from
the negative real direction. The finitely many remaining $A_j$
have finitely many possible nonzero lattice differences $A_j-A$.
Equation~\eqref{eq:mixed-small-denominator} consequently gives
finite-order poles at finitely many positive-multiple lattices in
$C_\bullet$, with fixed additional powers of $|C_\bullet|$.
For clarity, the fixed residue jets can be estimated directly,
without a fixed-radius circle near a real pole. If the unequal-$A$
denominators at the residue are $d_1,\ldots,d_s$, then
\[
 \left.\frac{\partial^q}{\partial t^q}
       \prod_{\ell=1}^s(t+d_\ell)^{-1}\right|_{t=0}
 =(-1)^q q!\left(\prod_{\ell=1}^s d_\ell^{-1}\right)
   \sum_{q_1+\cdots+q_s=q}\prod_{\ell=1}^s d_\ell^{-q_\ell}.
\]
Here $q$ is bounded in terms of the fixed main charge $A$.
Only boundedly many $d_\ell$ can be small, by the sorted-cost
estimate. The formula adds at most $q$ inverse-distance powers
and a factor bounded by a fixed multiple of $(1+r)^q$.
The latter is at most $C_q2^r$ and so can be absorbed in a single
larger geometric base, independent of the charge budget. Thus
the coefficient bound is $K^{(A)}K_*^{r+1}/r!$ times the stated
finite pole and polynomial losses. Equivalently a Cauchy circle
could be taken with radius at most half the smallest $|d_\ell|$;
that radius is not fixed near a lattice pole.

\subsubsection{Control of the meromorphic coefficients}

The following lemma controls these lattice poles throughout the later
bands, including the real columns required by the packet definition.
It will also apply to the one-sided charts in the fast-damping case.

\begin{lemma}[Meromorphic losses and their derivatives]\label{lem:mixed-poles}
Suppose a fixed outer formula has only finite-order poles on
lattices of finitely many positive multiples of $C_\bullet$, with
fixed polynomial factors in $C_\bullet$ and retained primary inputs,
and its remaining analytic factors are uniformly bounded on the
smaller argument boxes.
It then has the raw growth and simultaneous safe-column improvements
of Definition~\ref{def:packet}, with every prescribed finite number
of independent derivatives. This remains true with a scalar
detuning $\beta_s$, analytic with the same smaller-box margins, satisfying
$|\beta_s|\leq K e^{-\omega\Re R}$ before the first of the $n,k$
transitions. At a transition $j\leq k$, unit arguments in
$C_\bullet$ may vary on relative radii
\begin{equation}\label{eq:mixed-unit-radius}
 \eps\min(1,\Lambda_k/\Lambda_j).
\end{equation}
On a general usable $j$ band replace $\Lambda_j$ by
$\Lambda_{j-1}$; on safe sets a sufficiently small inverse power of
$|Z_k|$ also suffices.
\end{lemma}

\begin{proof}
For real $a\ne0,b$,
\[
 |a+b/C_\bullet|
 =|a|\frac{|C_\bullet+b/a|}{|C_\bullet|}
 \geq |a|\frac{|\Im C_\bullet|}{|C_\bullet|}.
\]
The same calculation applies to a fixed multiple of the ratio and
to any fixed pole order. Equations~\eqref{eq:mixed-ratio-im} and the
entering angular scale give logarithmic loss
$O(f_j(u)+\log(2+\Lambda_{j-1}))$. The flag estimate
$\log(2+\Lambda_{j-1})\leq C D_{j-1}+o(\Re Z_j)$ places this within
the permitted raw growth. At the far side $|v_{\rm ang}|\asymp
1/\Lambda_j$, and the same loss is $O(f_j+\log(2+\Lambda_j))$;
there is no exceptional constant there.

For completeness, verify the hypotheses of
Lemma~\ref{lem:flags-safe-columns} for the actual pole sweeps. A fixed
positive multiple has real value $x=\exp(f_k+\gamma)$ with
$\gamma',\gamma''=O(1)$: the bounded backgrounds have bounded real
derivatives and the differentiated strict-small unit corrections
tend to zero. Hence
\begin{equation}\label{eq:mixed-sweep}
 \frac{x'}x=\Lambda_k+O(1),\qquad
 \frac{x''}{(x')^2}
 =\frac1x\left(1+
  \frac{\Lambda_k'+\gamma''}{(\Lambda_k+\gamma')^2}\right)
 =\frac{1+o(1)}x.
\end{equation}
The inverse speed varies by $1+o(1)$ over a fixed lattice period.
Thus a small fixed lattice collar occupies an arbitrarily small
proportion of each fixed path interval; incomplete end periods have
vanishing length. For a finite list the sum of the bad proportions
is less than one. At the remaining columns real detuning persists
at small imaginary parts by Equation~\eqref{eq:mixed-real-shift};
at larger imaginary parts the imaginary separation suffices.

In the exceptional entering collar put $D=D_{j-1}$ and
$F=f_j(u_*)$. If $D>e^F/F^{K+3}$, then $F=o(D)$, while the uncapped
term defining $D$ gives $\Lambda_{j-1}\gtrsim e^D$. The window
$|\Delta f_{j-1}|\leq C_1\log(2+D)$ has path width
$O(e^{-D}\log(2+D))$. Within it $f_j\leq F+O(\log D)=o(D)$ and
the sweeps with $k\geq j$ satisfy
$x'=O(Z_j\Lambda_j)=e^{o(D)}$. Their changes are therefore $o(1)$.
This verifies the required persistence on the entire exceptional
collar, with one choice for a finite list.

Relative variations of the units on
Equation~\eqref{eq:mixed-unit-radius} change $C_\bullet$ by at most
a small fixed fraction of its imaginary separation. The actual
first-$j$ small arguments are exponentially smaller than these
radii. At safe points inverse-power radii preserve real detuning.
In the fast-damping construction below the balanced detuning has
size $O(e^{-\omega\Re R})$. Before $\min(n,k)$ is passed,
\[
 \Lambda_{j-1}\leq\Lambda_i
 =\Lambda_n+\Lambda_k+O(1),\qquad
 |\Im C_\bullet|\gtrsim
 |C_\bullet|\frac{\Lambda_k}{\Lambda_n+\Lambda_k+O(1)}.
\]
The logarithm of the inverse lower bound is
$O(\log(2+\Lambda_n))=o(\Re R)$. The detuning cannot remove the
separation. On safe real columns it is likewise eventually smaller
than the chosen lattice gap.

Take frozen holomorphic choices first. The finite triangular
formulas permit local independent-variable Cauchy radii
$\exp(-C f_j(u))$, multiplied when necessary by the pole-preserving
factor just obtained. Fixed differentiation increases only fixed
pole powers and these polynomial losses. On both sides of a $j$
transition the loss has logarithm $O(f_j)$, after taking extra
angular room; choose deeper undifferentiated accuracy first.
The general-band additional loss is precisely the allowed
$O(\log(2+\Lambda_{j-1}))$. Derivatives of the smooth cutoff choices
have the same bounds. These arguments apply on the small tubes of
controlled high real jets; they do not require a fixed ordinary
complex disk around a rotating dependent scale.
\end{proof}

\subsubsection{The leading transition and cutoff retirement}

We return to the retained low-$A$ residue sum. At the leading fringe,
all its pole losses and the fixed powers of $L,C$ have logarithm
$o(\Re L)$, by Lemmas~\ref{lem:mixed-ratio} and~\ref{lem:mixed-poles}.

For a requested depth $N$, include all $A\leq6N$ with these
cutoffs. Lemma~\ref{lem:mixed-tree} shows that this includes every
projected-low pole; the additional included poles have
$A+\mu B\geq N'$ and are individually exponentially small at
that depth after its fixed losses. There are only boundedly many
such low-$A$ indices per simplex. Thus these expressions, summed
against the analytic weights, give the leading transition.
One may move its finite budget endpoints to avoid charge values.

Choose a positive lattice unit $\omega$ for the $B$ values and put
\begin{equation}\label{eq:mixed-small-D}
 D=e^{-\omega W}.
\end{equation}
Since $B\geq-K_0A$, at fixed $A$ one can first extract a fixed
negative power of $D$. The remaining sum uses nonnegative integer
powers of $D$. The inequality $|B|\leq K_0M_0$ makes it normally
convergent on a fixed small $D$ disk, and even on a somewhat larger
disk if the original weights are further reduced. The powers
$L^dC^t$ are exact monomial prefactors times units. In particular,
once $e^{-AL-BW}$ has been split using the affine primary inputs,
it is never recomputed by varying the artificial ratio arguments
of a coefficient.

\begin{lemma}[Small-damping formal replacement]\label{lem:mixed-small-formal}
Before the first of the $n,k$ transitions, changing the moving
cutoff in a fixed coefficient changes its actual value, with
fixed derivatives, by
\begin{equation}\label{eq:mixed-mixed-cost}
 O\left(\exp(-a|Z_k|\Re R)\right)
 =O\left(\exp\{-a' e^{\Re f_i}\cos(\Im f_n)\}\right).
\end{equation}
If $n<k$, a fixed Taylor degree in $D$ removes that cutoff at the
$n$ transition. At the $k$ fringe the remaining coefficients have
Gevrey--$1$ expansions in $h=1/C_\bullet$, with optimal error
$O(e^{-a|Z_k|})$, uniformly on their smaller analytic argument
boxes. If $k\leq n$, the same assertion holds before $D$ is
dispersed, after replacing the cutoff sum by its full sum.
\end{lemma}

\begin{proof}
Neighboring cutoffs differ only in positive $B\geq a|Z_k|$.
The factors $e^{-BW}$ give Equation~\eqref{eq:mixed-mixed-cost};
the finite pole and prefactor losses are absorbed by decreasing
$a$. Equation~\eqref{eq:mixed-log-tie} identifies its two costs.
If $D$ is dispersed first, any fixed coefficient uses a fixed
finite set of $B$ values, so eventually it is independent of the
cutoff. Its unequal-$A$ factors retain the pole estimates already
proved.

At the $k$ fringe the ratio is transverse to the real axis. For
real lattice $a\ne0,b$,
$|a+b/C_\bullet|\gtrsim|a|$. Thus all $B$ can be included, with
the same factorial gain; the omitted terms have $M_0\gtrsim|Z_k|$
and cost $e^{-a|Z_k|}$ in the geometric norm, also on the chosen
$D$ disks. For an individual tree, after the explicit powers
$C^t$ have been extracted, the functions in
Equation~\eqref{eq:mixed-small-denominator} are analytic for
\[
 |h|<\frac{\eps_A}{1+M_0},
\]
with bound $K^{(A)}K_*^{r+1}/r!$: the unequal-$A$ lattice gap and
$|B_j|+|B|=O(M_0)$ prove this directly. Its $h^q$ coefficient
is bounded by the same prefactor times $K_A^q(1+M_0)^q$.
For $0<\eta<1$,
\[
 \sum_{M\geq0}\eta^M(1+M)^q\leq K^{q+1}q!;
\]
for example compare the sum with the integral of
$e^{-a x}(2+x)^q$ and expand the latter polynomial. Summing the
tree bounds proves the Gevrey estimate.

Set $Q=|C_\bullet|$ and truncate at $J=\lfloor\eps'Q\rfloor$.
For $M_0\leq\eps''Q$, the coefficient disk is larger than the
used $h$ by a fixed factor; its Taylor remainder is exponentially
small in $J$. For $M_0>\eps''Q$, the actual terms still have their
uniform imaginary-sector bounds. The truncated terms lose at most
\[
 \bigl(1+K(1+M_0)/Q\bigr)^J.
\]
Choose $\eps'$ small: its logarithm is then at most a fixed small
multiple of $M_0$, and the geometric weight still leaves
$e^{-aM_0}$. Summing this tail proves the optimal remainder.
The proof on fixed analytic disks also gives the fixed jets by
Cauchy estimates. When $D$ was dispersed first it applies directly
to each of its fixed coefficients.
\end{proof}

\subsection{Fast damping: one-sided charts and a common normalization}

Now $P=W$, $R=L$, and $C=c=W/L\to\infty$. Write the layer
charges as $A=a\cdot p$, $B=b\cdot q$ in a common positive
rational lattice $\omega\mathbb N$. The layer derivation
\[
 \Delta=\sum_\ell(-a_\ell X_\ell)\partial_{X_\ell}
              +\sum_\ell b_\ell Y_\ell\partial_{Y_\ell}
\]
has eigenvalue $B-A$ on $X^pY^q$; let $\Pi$ project onto $A=B$.
We use different normalizations at the two endpoints. The output
graph equation has divisors $c+B-A$, so a bound on $A$ protects its
inverse while leaving the right charges unrestricted. At state degree
$m\geq1$, the equation for the logarithm of the input-coordinate derivative
has divisors $B-A-mc$, so its inverse is protected by a bound on $B$.
Later coordinate inversions may still have lattice poles, controlled
by Lemma~\ref{lem:mixed-poles}. Their common charge range will identify
the two coordinates at the midpoint, where both omitted endpoint layers
are small.

Take $N$ a small fixed multiple of $|Z_k|$, locally frozen, so
$N\leq\eps|c|$ even on the required variant neighborhoods.
The output, input, and joint rings retain respectively $A<N$,
$B<N$, and both inequalities. They are quotient rings by monomial
ideals. Equip them with absolute coefficient norms on fixed small
layer polydisks, and on a state disk strictly inside a larger one.
The field norms, including the finite margins needed for its
derivatives, are small because $g(0,0,v)=0$.

\begin{lemma}[Compatible layer charts]\label{lem:mixed-charts}
In these rings there are an output graph $\phi$, an input
coordinate $\Phi$, an output coordinate
$T=\phi+\sum_{m=1}^Mt_m u^m$ for each fixed $M$, and a balanced
polynomial $\beta$ without constant term. Their constructions are
compatible under smaller charge quotients. Keeping a scalar
$\beta_s$ external until the final substitution, their only
remaining poles are finite-order poles on finitely many
positive-multiple lattices of $c-\beta_s$. In the joint ring,
after $\beta_s=\beta$,
\begin{equation}\label{eq:mixed-chart-identities}
 \Phi(T(u))=u\pmod{u^{M+1}},\qquad
 u'=(-c+\beta)u.
\end{equation}
\end{lemma}

\begin{proof}
In the output ring solve
\begin{equation}\label{eq:mixed-graph}
 (\Delta+c)\phi=g(X,Y,\phi).
\end{equation}
Here $B-A\geq-N$, and the allowed argument of $c$ gives
$|B-A+c|\geq a|c|$; thus the inverse has norm $O(1/|c|)$.
The small field norm makes Equation~\eqref{eq:mixed-graph} a
contraction in a fixed small graph ball. Put
\begin{equation}\label{eq:mixed-output-linear}
 b=g_v(X,Y,\phi),\quad \beta_o=\Pi b,\quad
 t_1^0=\exp\bigl(\Delta^{-1}(b-\beta_o)\bigr),
\end{equation}
where the inverse is zero on balance. The rational lattice gives
a bounded off-balance inverse, so $t_1^0$ is a bounded unit close
to one. Its eventual balanced normalization is specified below.

In the input ring solve
\begin{equation}\label{eq:mixed-input-H}
 (\Delta-cv\partial_v)H_1
       =\beta_e-g_v-g\partial_v H_1,
 \qquad \Pi[H_1]_{v^0}=0.
\end{equation}
Choose $\beta_e$ to cancel the omitted balanced $v^0$ projection
of the right side. At degree $v^m$, $m\geq1$, the inverse is
$O((m|c|)^{-1})$, since $B-A\leq N$. In a product
$g_pv^p\partial_v(H_l v^l)$ of degree $m$ one has
$l=m-p+1\leq m+1$, so the factor $l$ from differentiation divided
by $m$ is at most two. At degree zero only $l=1$ can contribute,
and off balance the inverse is bounded by the lattice gap.
In the absolute coefficient norm these bounds make the operator
on $H_1$ contractive when the $g$ norm is small. There is no
shrinking state radius in successive iterations: the derivative
has already been compensated coefficientwise. Thus $H_1$ is
small, $\Psi=e^{H_1}$ is a bounded unit close to one, and
$\beta_e$ is a small balanced polynomial without constant term.

Introduce the external scalar $\beta_s$ and define
\begin{equation}\label{eq:mixed-input-Phi}
 \Phi=f+\int_0^v\Psi(X,Y,t)\,\dd t,
 \qquad (\Delta+c-\beta_s)f=-g(X,Y,0)\Psi(X,Y,0).
\end{equation}
The layer-constant term of $f$ is zero. This final inverse is
meromorphic; it is not used in solving
Equation~\eqref{eq:mixed-input-H}.

In the joint ring, Equation~\eqref{eq:mixed-graph} and the chain
rule in Equation~\eqref{eq:mixed-input-H} give
\[
 \Delta(H_1(X,Y,\phi))=\beta_e-b.
\]
Balanced projection yields $\beta_e=\Pi b=\beta_o=:\beta$.
These are equal as polynomials, since balance in either one-sided
ring already lies in the joint ring. Moreover
$\Delta\log(\Psi(\phi)t_1^0)=0$. Lift the balanced reciprocal
of $\Psi(\phi)t_1^0$ from the joint ring to the output ring and
multiply $t_1^0$ by it. The resulting bounded unit $t_1$ satisfies
\begin{equation}\label{eq:mixed-normalization}
 \Psi(\phi)t_1=1
\end{equation}
in the joint ring. Taking this balanced reciprocal and its lift
commutes with every smaller charge quotient.

For $m\geq2$ put $t_m=t_1s_m$ and solve successively
\begin{equation}\label{eq:mixed-sm}
 [\Delta-(m-1)(c-\beta_s)]s_m
  =t_1^{-1}[g(X,Y,T)-g(X,Y,\phi)-b(T-\phi)]_{u^m}.
\end{equation}
The right side uses only lower fiber orders. Keeping $\beta_s$
scalar makes this one diagonal inversion at each step. For fixed
$M$ only finitely many pole factors and pole powers occur; they
are at lattice values of the multiples $(m-1)(c-\beta_s)$.
The input inverse has the same property. In the joint ring
$|B-A|<N$, so all these inverses are uniformly safe for $\beta_s$
in a fixed scalar disk containing $\beta$ with a strict margin.

Substitute $\beta_s=\beta$ in the joint ring. To check the input
conjugacy explicitly, set
\[
 F=\Delta\Phi+(-cv+g)\Phi_v-(-c+\beta)\Phi.
\]
Equation~\eqref{eq:mixed-input-H} gives $F_v=0$, and
Equation~\eqref{eq:mixed-input-Phi} gives $F(0)=0$. Thus $F=0$.
The graph equation, $\Delta t_1=(b-\beta)t_1$, and
Equation~\eqref{eq:mixed-sm} prove the opposite conjugacy through
degree $M$. If $K(u)=\Phi(T(u))$, these two identities make its
constant term annihilated by $\Delta+c-\beta$ and its degree-$m$
coefficient, $m\geq2$, annihilated by
$\Delta-(m-1)(c-\beta)$. Those joint inverses are safe. The
constant term is therefore zero, the linear term is one by
Equation~\eqref{eq:mixed-normalization}, and all degrees two
through $M$ vanish. This proves
Equation~\eqref{eq:mixed-chart-identities}.

All equations, projections, products, and state substitutions
preserve the charge ideals. Uniqueness of the contractions and
the diagonal inversions proves compatibility under restriction,
first with $\beta_s$ held scalar and then under its fixed jets or
formal substitution. This also proves the asserted uniform joint
bounds and the description of the one-sided poles.
\end{proof}

\begin{lemma}[Fast-damping leading transition]\label{lem:mixed-fast-transition}
On the $i$ fringe, for every fixed desired depth in $\Re W$ the
passage has the finite expansion
\begin{equation}\label{eq:mixed-fast-expansion}
 \phi_{\rm out}+\sum_{m=1}^{M}
 e^{-mW}e^{m\beta_sL}t_{m,\rm out}\Phi_{\rm in}^{m},
 \qquad \beta_s=\beta(X,Y),
\end{equation}
up to that depth, with $M$ sufficiently large. Each fixed-$m$
coefficient is independent of how much larger a comparison order
is chosen.
\end{lemma}

\begin{proof}
Use the input coordinate on the first half of the passage, compare the
two coordinates in their joint charge range at the midpoint, and then
use the output coordinate on the second half. Damping makes the fiber
variable small by the midpoint, while the omitted endpoint charges are
small on their respective halves.

Along actual profiles, a balanced monomial is constant and has
size $O(e^{-\omega\Re L})$. Consequently $\beta_s$ is constant,
$|\beta_s|\leq K e^{-\omega\Re L}$, and $|L\beta_s|\ll1$.
Equation~\eqref{eq:mixed-cos-gap} gives
$|\Im c|\gtrsim |c|\Lambda_k/\Lambda_i$, and this is much larger
than $|\beta_s|$ by Lemma~\ref{lem:mixed-poles}. Put
\[
 \mathcal P=K_0\left(1+|\Im(c-\beta_s)|^{-1}\right).
\]
At each fixed fiber order the one-sided coefficients and the
needed state or layer derivatives have bounds by fixed powers of
$\mathcal P$ and $1+|c|$. Here
$\log\mathcal P=O(\log(2+\Lambda_n))=o(\Re L)$.

Split a monotone-action path at its midpoint $L/2$. On the first
half all omitted right charges have $B\geq N$ and pay
$e^{-N\Re L/2}$. Evaluating the input equations along the actual
solution gives
\[
 \frac{\dd\Phi}{\dd\theta}
       =(-c+\beta_s)\Phi+\epsilon_{\rm in},\qquad
 |\epsilon_{\rm in}|\leq
 K_M(1+|c|)\mathcal P^{K_M}e^{-N\Re L/2}.
\]
To justify this estimate even for the state derivative equation,
first use its polynomial balanced $\beta$ in the quotient; its
actual evaluation is exactly $\beta_s$. The difference between
products taken before or after projection has only omitted right
charges. The constant equation uses the scalar directly.
Stable variation of constants shows that at the midpoint $\Phi$
differs from
\[
 u(\theta)=e^{(-c+\beta_s)\theta}\Phi_{\rm in}
\]
by the displayed bound times $O(|L|)$. On the second half,
$|u|\leq K\mathcal P^K e^{-\Re W/2}$.

At the midpoint restriction of both one-sided data to the joint
ring changes evaluation by at most
$K_M\mathcal P^{K_M}e^{-N\Re L/2}$. This holds for projected
products and compositions as well. The joint coefficients are
analytic on a fixed scalar $\beta_s$ disk, so substituting the
smaller balanced series and then projecting has the same error
as projecting before substitution. We can therefore use
Equation~\eqref{eq:mixed-chart-identities} at the midpoint.
The unused Taylor orders cost
$K_M\mathcal P^{K_M}|u|^{M+1}$; an inverse-power
$\mathcal P$ radius in $u$ keeps $T$ in its state disk. Since
$\Phi_v=\Psi$ is uniformly close to one, the actual midpoint
state and $T(u)$ differ by the sum of these errors and the
preceding path-length loss.

On the second half the approximate solution $T(X,Y,u(\theta))$
satisfies the original equation through degree $M$, up to omitted
left charges, again of size at most
$K_M(1+|c|)\mathcal P^{K_M}e^{-N\Re L/2}$. Higher degrees cost
$K_M\mathcal P^{K_M}|u|^{M+1}$. Comparison with the exact stable
propagator and the small integral of $|g_v|$ loses only a constant
and path-length factors. At this fringe
\[
 \frac{N\Re L}{\Re W}\asymp
 \frac{\cos\arg L}{\cos\arg W}\longrightarrow\infty.
\]
All polynomial factors have logarithm $o(\Re W)$. Thus the charge
errors beat every fixed depth, and choosing the comparison order
$M$ beyond twice that depth handles the Taylor error. Evaluating
$T$ at the output gives Equation~\eqref{eq:mixed-fast-expansion}.
The triangular recursion for a fixed $t_m$ does not involve the
later comparison orders.
\end{proof}

\subsection{Fast-damping dispersion and the full charts}

Put $D=e^{-\omega L}$. At the input a right charge contributes
$D^{B/\omega}$, and at the output a left charge contributes
$D^{A/\omega}$. The balanced construction is pole-free and gives
\begin{equation}\label{eq:mixed-beta-D}
 \beta_s=D\widetilde\beta(D)
\end{equation}
with $\widetilde\beta$ uniformly bounded and analytic on small
$D$ disks. Treat $LD$ as a separate strict-small argument in
$e^{m\beta_sL}$; it is a finite sum of monomials times ordinary
factors through the primary log formulas, with the same first
index $n$. This prevents the unbounded $L$ from being used as an
ordinary parameter of that exponential.

Before $\min(n,k)$ is passed, the actual $D$ is much smaller than
the imaginary separation of $C_\bullet$. At an $n$ transition
with $n<k$, a disk
\begin{equation}\label{eq:mixed-D-radius}
 |D|\leq\eps\min(1,\Lambda_n^{-1})
\end{equation}
still preserves that separation: by
Equation~\eqref{eq:mixed-beta-D} its detuning is $O(\Lambda_n^{-1})$,
whereas $|C_\bullet|\Lambda_k\to\infty$. Its Taylor tail of
degree $J$ is bounded, up to fixed losses, by
\[
 \exp(-J\omega\Re L+O(J\log(2+\Lambda_n))),
\]
which supplies arbitrary fixed $n$ depth. The denominator jets
after this expansion use $\beta_s=0$, leaving finite powers of
the original lattice poles.

Every fixed $D$ degree uses only bounded output left charges and
bounded input right charges. Lemma~\ref{lem:mixed-charts} proves
that these coefficients are independent of the growing cutoff:
restriction first at scalar detuning, balanced normalization on
the joint intersection, and finite scalar differentiation all
commute with the smaller charge quotients. Before this first
dispersal, neighboring cutoffs differ at actual endpoints only
by omitted charges or balanced terms of size
$e^{-aN\Re L}$. The change in scalar detuning has the same size;
finite pole powers and derivatives preserve it after reducing
$a$. This is exactly Equation~\eqref{eq:mixed-mixed-cost}, now
with $R=L$. In comparing ratio variants keep $D$ independently
small on its allowed disk; the already extracted shifts stay
fixed.

At the $k$ fringe use the full convergent layer rings, with no
charge quotients. On a fixed small ratio-unit disk, transversality
gives, for real layer eigenvalues $d$ and integers $l\geq1$,
\begin{equation}\label{eq:mixed-transversal}
 |d\pm lC_\bullet|\gtrsim |d|+l|C_\bullet|.
\end{equation}
All constructions in Lemma~\ref{lem:mixed-charts} now work in
these full rings; keep $\beta_s$ external on a bounded scalar
disk and keep $D,LD$ on separate small disks. The same balanced
normalization and conjugacy proof applies. If $k\leq n$, the
partial charts are their compatible projections. On slightly
enlarged amplitude and $D$ boxes their omitted endpoint charges
cost $e^{-aN}=e^{-a'|Z_k|}$, also after factoring $D$ from
$\beta_s$. The bounded analytic substitutions preserve that
precision. If $n<k$, the previously obtained fixed $D,LD$ jets
agree directly with those of the full charts.

\begin{lemma}[Gevrey estimates for full charts]\label{lem:mixed-fast-gevrey}
For each fixed output fiber order, the full charts and their
bounded analytic endpoint operations have consistent Gevrey--$1$
formal expansions in $h=1/C_\bullet$. On the transverse $k$
fringe they agree with sufficiently small proportional optimal
truncations to precision $O(e^{-a|Z_k|})$. The assertion is uniform
on smaller disks in all the ordinary and small arguments, and
holds for their fixed jets.
\end{lemma}

\begin{proof}
Multiply the graph equation by $h$; its linear part is
$1+h\Delta$. At state degree $m\geq1$, divide the input equation
by $-cm$. Its leading inverse is regular and its perturbations
are $h\Delta/m$ and the small product operator
$h[g\partial_vH_1]_m/m$. At $m=0$ retain the off-balance
$\Delta^{-1}$ equation. Only degree one of $H_1$ can enter that
last differentiated product. The $f$ and fixed-$s_m$ equations,
after multiplication by $h$, likewise have regular leading
inverses.

Here are uniform coefficient estimates for these formal
equations. For a final budget $J$ use absolute formal weights
$(\eps/(J+1))^q$ at order $h^q$, and layer radii whose logarithms
decrease linearly with $q/(J+1)$ between two fixed nested disks.
Products respect the associated coefficient norm. Applying
$\Delta$ while advancing the formal degree once costs at most
$K(J+1)$ by the one-step radius loss, which the $h$ weight
reduces to $K\eps$. Move those terms to the right of the leading
inversion. The graph is a contraction for small $\eps$ and small
$g$. For $H_1$ the state derivative is still compensated by
$l/m\leq2$; the $m=0$ coupling is small at the same fixed state
radius and has a bounded off-balance inverse. These estimates
give a uniform contraction on the entire finite formal norm,
and hence a bound independent of $J$. The equations determine
the same coefficient at order $q$ for every $J\geq q$.
Taking $J=q$ gives $K^{q+1}(q+1)^q$, and hence
$K_1^{q+1}q!$, for that coefficient.

Balanced projection and the off-balance primitive are bounded
in these norms. Analytic operations on units with a fixed
nonzero limiting value preserve them on smaller disks; this is
also the single-variable case of
Lemma~\ref{lem:additive-gevrey-operations}. The balanced
normalization and the finite triangular $f,t_m$ operations
therefore have the same factorial estimates.

For the actual comparison take $J=\lfloor\eps'|C_\bullet|\rfloor$
with $\eps'$ small enough that the formal dummy radius is at
least twice the used $|h|$. Formal cancellation and the bounded
norm just proved make the residual of the truncated equations
$O(e^{-a|C_\bullet|})$ in their scaled norms on smaller layer
disks; one extra fixed disk shrink pays a final $\Delta$.
Equation~\eqref{eq:mixed-transversal} bounds the actual scaled
inverses. The same small Lipschitz constants compare the actual
graph and input fixed point with the approximations, preserving
the residual precision. The remaining bounded analytic and
triangular operations preserve it in turn. Endpoint substitution,
factoring $D$ from $\beta$, and the scalar substitution take place
on strictly smaller fixed argument boxes and are bounded
analytic operations. Cauchy estimates there prove the same
assertions for fixed jets. This includes the jets already used
when $D,LD$ were dispersed before reaching this fringe.
\end{proof}

\subsection{Dispersal and source retirement}

The preceding lemmas give the actual leading band and its
arbitrary-depth expansion. We now construct the later packet bands.
The order of the two unbalanced replacements determines when a
coefficient ceases to depend on a discarded free input.

First extract every explicit affine exponential and every
monomial prefactor. At each subsequent index Taylor-expand the
strict-small arguments whose first index is current, keeping the
other arguments as variables and setting the dispersed arguments
to zero in coefficient evaluations. Fixed derivatives of already
dispersed coefficients obey the same subsequent estimates:
their formulas are the corresponding fixed jets, and their
Cauchy margins were reserved before that dispersal. In the unbalanced
cases, before the ratio's $k$ transition these evaluations use
$C_\bullet$ and the radii in Lemma~\ref{lem:mixed-poles}. At the $k$
fringe fixed ratio disks are available.

In either unbalanced case the moving charge cutoff exists only
on bands $i<j\leq\min(n,k)$. If $n<k$, disperse $D$ (and $LD$
when present) first. Every fixed coefficient then has bounded
charges and no numerical cutoff. If $k\leq n$, first pass to the
full sum or full charts and their formal expansions at $k$;
disperse $D,LD$ when their first index is reached. If $n=k$, the
full/formal replacement is done on the small $D,LD$ disks before
the Taylor dispersal at that same transition.

After the formal replacement write, by
Equation~\eqref{eq:mixed-log-tie},
\[
 h=C_\bullet^{-1}
  =\frac{A_R}{A_P}e^{-b_*}e^{-f_k}
       \frac{1+\Theta_R^{\bullet}}{1+\Theta_P^{\bullet}}.
\]
Absorb its bounded unit into the formal coefficients, preserving
their Gevrey bounds. The monomial $e^{-f_k}$ has first index
$\sigma(k)>k$. Retain a small proportional optimal polynomial
until that index. Its tail from any fixed degree $q$ is bounded
by $K_q|e^{-f_k}|^q$, since its terms decrease geometrically up
to the chosen cutoff. Thus only fixed sufficiently many degrees
remain at the $\sigma(k)$ transition. Beyond it no expression
uses the numerical value of that cutoff. These are precisely the
finite-argument dispersal operations justified in
Proposition~\ref{prop:additive-dispersal}; the estimates here
also prove them directly for the additional meromorphic stage.

For clarity, verify every source cost. Before $\min(n,k)$ has
passed, neighboring cutoffs have cost
\[
 U=e^{\Re f_i}\cos(\Im f_n)\asymp |Z_k|\Re Z_n.
\]
It is admissible because $i<j$, $\sigma(i)>\min(n,k)\geq j$,
and $n\geq j$. Throughout its active band,
\begin{equation}\label{eq:mixed-cost-domination}
 \frac{U}{\Re Z_j}\gtrsim
 e^{f_i(u)-f_j(u)}
   \frac{\cos(\Im f_n)}{\cos(\Im f_j)}
 \gtrsim e^{f_i(u)-f_j(u)}\longrightarrow\infty.
\end{equation}
This remains true on the thin side. In addition
$\Lambda_{j-1}\leq\Lambda_i=\Lambda_n+\Lambda_k+O(1)
\leq2\Lambda_j+O(1)$ there. An exceptional uncapped constant,
if present, is consequently only
$D_{j-1}=O(\log(2+\Lambda_j))$ in these active bands. All raw,
normalizing, pole, and fixed derivative losses therefore have
logarithm $o(U)$. At the owning $i$ fringe
$U/\Re Z_i\asymp\cos\arg R/\cos\arg P\to\infty$,
so its error also pays arbitrary leading transition depth.

After full/formal replacement the optimal-polynomial cost is
the simple cost $e^{\Re f_k}$; it is used only on bands
$k<j\leq\sigma(k)$. The simple source estimate from the flag
calculus absorbs all fixed losses there. At the replacement
itself $e^{-a|Z_k|}$ pays every fixed $k$ fringe depth, since
$\Re Z_k/|Z_k|\to0$. Thus neither type of source survives the
last index at which its defining logarithm is retained.

Choose frozen orders and cutoffs on fixed-factor bins of
$e^{\Re f_k}$ and interpolate smoothly on overlaps. The
differences just estimated have the appropriate exponential
precision; the partitions and their fixed derivatives have
logarithmic losses negligible relative to those costs. These
choices use only retained formulas and can be made simultaneously
for every finite request. There are no sources in band one.
The actual passage through the leading band is genuinely
holomorphic by Lemma~\ref{lem:mixed-actual}.

\subsection{Controlled inputs and terminal coefficient germs}

\begin{lemma}[Uniformity for one controlled family]
\label{lem:mixed-controlled-uniformity}
The mixed constructions and their finite dispersion operations have
the controlled-family uniformity of
Definition~\ref{def:calculus-controlled-family}. For one fixed finite
request, their constants and numerical far thresholds depend only
on its common analytic boxes and margins, field bounds, finite
rate and affine data, and finite quantitative flag and input-jet
controls. They do not depend on when an individual input path
returns to a smaller neighborhood of its ordinary germ.
\end{lemma}

\begin{proof}
Fix the finite labels, derivative orders, and Taylor and transition
budgets in the request. Choose common compact ordinary, layer, and
state boxes with positive distance to the larger analytic domains.
The field suprema and their finitely needed Cauchy bounds are then
common. The initial amplitude choice gives a fixed contraction
margin, say at most $1/2$, in the actual integral equations and
in the layer-ring fixed-point equations. For the latter the graph
inverse, the coefficientwise $l/m\leq2$ estimate, and the bounded
off-balance inverse give the same margin at every point satisfying
the prescribed numerical conditions on $|c|$ and the charge cutoff.
These estimates involve values in these boxes, not the subsequent
history of the ordinary test in the flag path parameter.

The angular comparisons use the common finite real jets and Taylor
remainders only through finite triangular sensitivity bounds.
Their thresholds can therefore be chosen as common inequalities
in the finitely many rates, gaps, and angular buffers. In the
comparable case first fix $\kappa$ and the side constant as in
Equation~\eqref{eq:mixed-comparable-detuning}; its vanishing
remainder has the common modulus provided by the family controls.
This gives common normalized action and pole-cone margins. The
polygon length constants and stable comparison bounds are then
common by Lemma~\ref{lem:mixed-actual}. A dependence on the fixed
limiting slope is permitted here; common amplitude smallness over
slopes is the stronger, separately proved assertion above.

For the trees, geometric weights, counting constants, and the
sorted-cost constant are determined by the field boxes and finite
rate list. For each of the finitely many requested main charges,
the finite low-cost separations and derivative orders give one
common constant; the tail has the common factorial bound of
Lemma~\ref{lem:mixed-tree}. Its contour remainders require only
specified inequalities such as
$\Re L>K\log(1/\rho)$ and the finitely many cap-separation
thresholds. The small-slope formal replacement uses the same
geometric moment estimate at every point. Similarly, the finite
formal norms in Lemma~\ref{lem:mixed-fast-gevrey} use two fixed
layer radii, a fixed state margin, and a common small formal
weight. Taking the minimum of the finitely many permissible
optimal proportions makes all factorial and optimal-error
constants common. An already fixed convention for a shorter
label is preserved by its exponentially small cutoff-comparison
estimate; this does not redefine that label on a new test.

On a meromorphic band, the finite lattice lists and maximal pole
orders are determined by the labels and their derivative orders.
The direct finite jet formula above gives common powers of the
explicit pole losses, even when the numerical denominators vary.
Lemma~\ref{lem:mixed-poles} retains the allowed $D_{j-1}$ loss on
the entire forward domain. Its safe-column proof is also uniform:
the common real-jet and flag controls make
Equation~\eqref{eq:mixed-sweep} uniform, so a fixed forbidden
lattice width and a fixed union bound leave columns in every
required fixed-length interval after common numerical thresholds.
The exceptional-collar width and sweep-speed estimates use the
same constants; no later return to the germ is used to obtain
their $o(1)$ variation.

Independent Cauchy radii are a common fixed exponential
$\exp(-C f_j)$ times the explicit pole-preserving factor. Their
finite powers have the common logarithmic bounds already proved.
Taylor depths can therefore be chosen once for the request,
before applying those losses. The moving-cutoff and optimal-order
comparisons have the common costs in
Equation~\eqref{eq:mixed-mixed-cost} and $e^{\Re f_k}$,
respectively. Their absorption uses the common numerical gap
and side inequalities in Equation~\eqref{eq:mixed-cost-domination}
and the simple-cost estimate. Fixed-width bins in the retained
log modulus have uniformly bounded overlap. The same partitions
in those bins have common derivative bounds after the stated
triangular losses; hence both the number and precision constants
of the source terms are common for this finite request.

Finally, high real-jet agreement preserves these domains, angles,
and evaluations of independent derivatives. The differentiated source
bounds come from the frozen holomorphic formulas and their proved
cutoff overlaps. Holomorphic ordinary inputs contribute no input defect;
smooth sheets already constructed by Theorem~\ref{thm:calculus}
supply their own differentiated defects.

For pullbacks, the finite chain-rule factors have exactly the raw,
Cauchy, and controlled-sheet losses just bounded. Their logarithms are
$o(U)$ at intrinsic active cutoff sources by the preceding cost
comparisons. Inherited sheet sources retain their own permitted costs;
the comparisons in Definition~\ref{def:packet} and
Equation~\eqref{eq:calculus-source-fringe} apply to them as well.
At a level-$p$ good join, keep a fixed
current loss as $C\Re Z_p$; the remaining later-rate and local Cauchy
losses are $o(V_p)$. Lemma~\ref{lem:calculus-source-pullback} therefore
preserves the source estimates, choosing a transition budget larger
than its current losses or a sufficiently vertical collar for an
independently established $e^{-cV_p}$ error. The common controls above
make these choices uniform for this finite request.

All the preceding estimates hold throughout the part of each full
forward domain on which the stipulated common ordinary boxes and
numerical inequalities hold. Entering a smaller germ neighborhood
may be delayed by a different amount on each path. It affects a
later request using that smaller box, not any constant, starting
inequality, or source count for this fixed request. This proves
the claimed uniformity.
\end{proof}

\begin{proof}[Proof of Theorem~\ref{thm:mixed-packets}]
We have constructed the actual passage, its transition expansions,
and the raw coefficients in every band. It remains to collect their
growth and derivative estimates and to verify support and germ
invariance.

The growth of a fixed raw label is at most
$\exp(C\Re Z_j+CD_{j-1})$: after explicit shift translations,
its bounded analytic pieces, factorially summed tree pieces, or
finite chart pieces have only the losses of
Lemma~\ref{lem:mixed-poles}. That lemma also supplies the side
and safe-column improvements, simultaneously for fixed
derivatives and the full exceptional collar. Its local Cauchy
estimates applied with deeper initial transition accuracy prove
all differentiated transition bounds. Intrinsic antiholomorphic
derivatives arise only in the smoothed primitive choices; their
differences and all fixed derivatives have exactly the active
precision costs. On previously constructed smooth input sheets
the inherited input sources are also propagated by the chain rule,
as proved in Lemma~\ref{lem:mixed-controlled-uniformity}.
These observations establish the required growth, transition, and
source estimates in the independent variables, including controlled
input tests.

Finally the support is iterated left-finite. At the first index
the main charges form a locally finite positive semigroup or
a positive lattice. At a fixed charge there are only finitely
many prefactor translations; all further operations take Taylor
powers of a finite set of strict-positive monomials, or the next
fixed lattice charges. Fixing each prefix bounds the degrees
and translations used at its next coordinate. Gevrey evaluation
is used only before its first index is extracted, so it introduces
no infinite set below a fixed owning budget. Colliding
coefficients are added by these finite recursive rules. At the
last band only ordinary analytic operations and normally
convergent bounded-argument sums remain.

The terminal recipes are invariant under an allowed saturated
refinement of the flag. In the comparable construction they are
the residues of the exact convolution formula, with collisions
interpreted by the analytic cluster integral
Equation~\eqref{eq:mixed-bounded-H}. Splitting a primary affine
shift into refined nodes changes only the representation of its
exponent, and regrouping a finite fixed-charge extraction leaves
these analytic residue identities unchanged. In the small-slope
construction each coefficient after cutoff retirement is a
coefficient of the same fixed residue product; its formal
$h$ coefficients are determined by that product, independently
of optimal order. In the fast-damping construction, uniqueness
of the graph and input contraction, the stipulated off-balance
inverse, the balanced normalization, and the finite triangular
fiber recursion determine the coefficient recipes. Compatibility
with smaller charge ideals shows that every fixed endpoint jet
is the same germ before and after removing its cutoff. The
full-ring formal recursions have exactly these jets. Thus all
terminal identities are identities of ordinary analytic germs,
with the same fixed lateral determinations.

The same argument also covers the small absolute perturbations of
old free inputs in Definition~\ref{def:elimination-retestable}.
Keep the selected affine flag formulas, rate lists, limiting slope
and residual regime, ordinary field germs, and lateral determinations.
Changing the numerical old free inputs within that regime changes
the evaluations of the extracted shifts and of the intermediate
coefficient formulas, but changes none of the residue identities,
charge projections, normalizations, or formal recursions just
listed. Their final coefficient formulas contain no numerical free
input or cutoff. They consequently give the identical ordinary
analytic germs, rather than merely the same values at the limiting
ordinary point. A change of modulus bin affects only an intermediate
cutoff choice, whose dependence has already disappeared by the
proved retirement rules. The estimates apply to the perturbed
sequence by the same controlled-test construction, with its
preserved regime gaps and smaller neighborhoods as needed.

An inserted unused node contributes an identity transition. Applying the earlier
finite ordinary-chart recipes then invokes
Lemma~\ref{lem:calculus-refinement}. This proves the stated retesting
and refinement invariance without identifying an unspecified flat error with
zero, and completes the theorem.
\end{proof}

\subsection{Uniformity of the primitive library}

\begin{proposition}[Quantitative primitive input bounds]
\label{prop:calculus-library-uniformity}
For the primitive constructions in Sections~\ref{sec:additive},
\ref{sec:regular}, and~\ref{sec:mixed}, the family requirement in
Assumption~\ref{ass:calculus-primitive} follows from the estimates on
common argument boxes proved in those sections.  For each fixed finite
request their constants and far-regime thresholds depend only on those
boxes and margins, the finite field and flag data, and the finite controls
in Definition~\ref{def:calculus-controlled-family}.  The time at which
an ordinary test returns to its germ does not enter these constants.
\end{proposition}

\begin{proof}
We specify the dependence in the estimates used there.  This also
distinguishes this assertion from merely assuming uniformity in addition
to estimates proved one test at a time.

For additive primitives, fix the larger field, profile, and state boxes
and their reserved margins.  The path integrals in
Equations~\eqref{eq:additive-slow-integral} and
\eqref{eq:additive-fast-integral} are bounded by constants times the
amplitude, with constants depending only on the fixed exponents and
boxes.  The kernel in Equation~\eqref{eq:additive-kernel} has the same
damping bound throughout those boxes.  For a fixed coefficient request,
its finitely many exceptional diagonals have common bounds; the infinite
opposite tail uses the common estimate
\[
 \|\mathcal L_m^{-1}\Pi_{>N}F\|
 \le (C_0/N)\|F\|.
\]
The threshold $N>2C_0\|A_0\|$ is selected once on the larger boxes,
as are the summable state and spatial margins.  In the formal estimate
the radius $\varepsilon/J$ pays a derivative of norm $CJ$; hence one
choice of $\varepsilon$ gives the Gevrey constants for that request.
The actual residual comparison, its fixed polynomial moments, and the
finite charge remainder then use these same constants.  The normalized
primitive integral in Equation~\eqref{eq:additive-inward-integral}
is uniformly integrable after choosing its fixed exponent margin.
The infinite anchor in that integral is in the independent spatial
variable $w$, while the ordinary arguments are held in their common
boxes.  It imposes no return-time requirement on an ordinary test
$t(s)$.  These observations apply to every estimate in
Lemmas~\ref{lem:additive-slow}, \ref{lem:additive-primitives}, and
\ref{lem:additive-gevrey-operations}, and hence to the finite formulas of
Propositions~\ref{prop:additive-charge} and
\ref{prop:additive-dispersal}.
Lemma~\ref{lem:additive-controlled-family} gives the resulting explicit
common jet and angular-error thresholds on the full forward domain.

For the one-rate primitives, the same common-box choice controls the
integrated layer norm and the stable variation-of-constants operator.
In Lemma~\ref{lem:regular-formal}, the nonzero charge gap is fixed and
the high-tail inverse again gains $C/N$.  The weighted norm for slow
degree has derivative cost $CJ$ multiplied by $\varepsilon/J$.
Consequently the pure and tail contraction margins, factorial constants,
optimal proportions, and the endpoint inverses in
Lemma~\ref{lem:regular-charges} are chosen on these boxes, before any
particular ordinary test is supplied.  The detailed controlled-input
verification is Lemma~\ref{lem:regular-controlled-inputs}.

For mixed primitives, the actual continuation in
Lemma~\ref{lem:mixed-actual} is controlled by an integrated fast-profile
norm and a small Lipschitz norm on the larger state box.  The contour
bounds of Lemma~\ref{lem:mixed-tree} use the fixed finite rate and charge
data.  Their infinite charge sums have a common geometric majorant
$\eta^M$ with $\eta<1$, and for a fixed derivative order $q$,
\[
 \sum_{M\ge0}\eta^M(1+M)^q\le C^{q+1}q!.
\]
For instance, with $a=-\log\eta>0$, the exponential-series inequality
$x^q\le q!\tau^{-q}e^{\tau x}$, applied with $\tau=a/2$, proves this
bound with a constant depending only on $\eta$.  Thus the fixed moments
and Gevrey constants of Lemmas~\ref{lem:mixed-small-formal} and
\ref{lem:mixed-fast-gevrey} are common.  The graph and layer contractions
in Lemma~\ref{lem:mixed-charts} have common diagonal inverse bounds on
the stipulated slope ranges.  The finite large-rate and detuning
inequalities in these arguments are precisely numerical regime
thresholds of Definition~\ref{def:calculus-controlled-family}; they
are independent of the history of the current inputs.
Lemma~\ref{lem:mixed-controlled-uniformity} records the complete uniform
choice through all of the mixed constructions.

The pole estimates also give uniform safe-column existence, rather than
requiring one numerical column for all tests.  To see this quantitatively,
let $x_a$ range over the finite list of positive real lattice sweeps in
Lemma~\ref{lem:mixed-poles}, and let $h_a>0$ be their lattice spacings.
Common first and second real-jet bounds give, beyond common flag
thresholds,
\[
 \frac{x_a'}{x_a}\ge\frac12\Lambda_{k(a)},\qquad
 \left|\frac{x_a''}{(x_a')^2}\right|\le\frac{C}{x_a}.
\]
On a complete lattice period where $x_a\ge X$, inverse speeds differ
by at most $e^{Ch_a/X}$.  A forbidden neighborhood of fixed width
$2\delta$ therefore occupies at most
$(4\delta/h_a)e^{Ch_a/X}$ of that period's traversal time.  The two
incomplete endpoint periods cost at most
$2h_a/\inf x_a'$.  First choose $\delta$ for the finite list, and then
choose common lower thresholds for $X$ and the speeds.  The union of
the forbidden parts occupies less than a prescribed fixed-length
interval, giving a safe column for each test with the same detuning
margin.  When the exceptional entering collar is needed, its change in
each sweep is bounded by
$C\exp(-D+o(D))\log(2+D)$, with a common modulus in the $o(D)$ from the
same finite jet and flag controls.  It is therefore smaller than the
chosen detuning margin.  This proves the uniform collar assertion while
leaving $D$ explicit.

Finally the independent-input Cauchy disks, frozen optimal orders,
and modulus-bin cutoffs in all three constructions are local in the
current retained tuple.  Their radii and loss exponents depend on
the common argument margins, finite pole orders, and quantitative flag
controls.  They are not disks of analytic continuation for the entire
test path.  Adjacent cutoff formulas have the stated exponential errors
on the larger common disks, so their fixed differentiated errors and
partition losses have common source constants.  A sufficiently high
fixed Taylor order preserves these geometric margins uniformly by
Lemma~\ref{lem:flags-background}; it does not itself create a source
estimate.  For source estimates under composition, the inputs are
holomorphic or have the already established differentiated defects of an
old chart sheet.  Every remaining factor in the finite chain rule has
one of the listed absorbable losses, uniformly for the controlled family.
Lemma~\ref{lem:calculus-source-pullback} therefore gives the composed
source bounds, including its separate treatment of current good-join
transition remainders and independently established join errors.

All these estimates are inequalities at the current arguments and hold
at every point of the full forward domain on which the common controls
hold.  Neither these norms nor the safe-column construction waits for an
individual test to reach a smaller ordinary neighborhood.  Such a wait
is used only for further preliminary requests in
Lemma~\ref{lem:calculus-uniform-separation}, whose constants disappear in
the final weighted estimate.  This proves the proposition.
\end{proof}
